\section*{Estado da correspondência}
As categorias abaixo reproduzem a matriz recebida. Correspondência inspecionada identifica declarações relevantes; integração à base exige as pontes documentadas. Contratos abstratos, cobertura parcial e pendências conservam seu alcance. O arquivo \texttt{mapas/MAPA\_TEX\_LEAN.md} reúne os nomes Lean, os caminhos e as posições de cada resultado.

Todos os 42 casos do dossiê de pendências estão fechados: 41 por prova Lean compilada e auditada e III.12.3 por prova textual sincronizada. Permanecem como ressalvas: a reverificação da leitura de $g(t_0)$ em I.4.3 contra o enunciado Lean; o enunciado de IV.5.4, cuja cota impressa era falsa para $M=0$ (contraexemplo documentado e certificado separadamente) e foi corrigida no texto; os modelos de custo de III.9.1 (operações unitárias; o desempate por $\kappa$ não está no programa denso) e de III.9.3 (o custo da aritmética racional em bits é hipótese do modelo). Os demais estados e ressalvas constam das linhas abaixo. As correspondências herdadas não são promovidas a provas verificadas por esta revisão.

Até 13/11/2026, 41 resultados antes pendentes passaram ao estado ``Formalizado e compilado'': provas Lean compiladas com Lean/Mathlib 4.28.0 nas bibliotecas \texttt{Pendencias} e \texttt{PendenciasParte1} dos ambientes \texttt{artigo\_2} e \texttt{base\_integrada}, sem \texttt{sorry} e apenas com os axiomas \texttt{propext}, \texttt{Classical.choice} e \texttt{Quot.sound}. Os demais estados e ressalvas seguem descritos acima.

\subsection*{Parte I}
\small
\begin{longtable}{@{}p{.11\linewidth}p{.09\linewidth}p{.09\linewidth}p{.59\linewidth}@{}}
\toprule Resultado & Original & Página & Estado e ressalva\\\midrule\endhead
\hyperref[map:artigo_1:2.1]{I.2.1} & 2.1 & \pageref{map:artigo_1:2.1} & Correspondência inspecionada\\[4pt]
\hyperref[map:artigo_1:2.2]{I.2.2} & 2.2 & \pageref{map:artigo_1:2.2} & Correspondência inspecionada\\[4pt]
\hyperref[map:artigo_1:2.3]{I.2.3} & 2.3 & \pageref{map:artigo_1:2.3} & Correspondência inspecionada\\[4pt]
\hyperref[map:artigo_1:2.4]{I.2.4} & 2.4 & \pageref{map:artigo_1:2.4} & Correspondência inspecionada\\[4pt]
\hyperref[map:artigo_1:2.5]{I.2.5} & 2.5 & \pageref{map:artigo_1:2.5} & Correspondência inspecionada\\[4pt]
\hyperref[map:artigo_1:3.1]{I.3.1} & 3.1 & \pageref{map:artigo_1:3.1} & Correspondência inspecionada\\[4pt]
\hyperref[map:artigo_1:3.2]{I.3.2} & 3.2 & \pageref{map:artigo_1:3.2} & Correspondência inspecionada\\[4pt]
\hyperref[map:artigo_1:3.3]{I.3.3} & 3.3 & \pageref{map:artigo_1:3.3} & Correspondência inspecionada\\[4pt]
\hyperref[map:artigo_1:3.4]{I.3.4} & 3.4 & \pageref{map:artigo_1:3.4} & Correspondência inspecionada\\[4pt]
\hyperref[map:artigo_1:3.5]{I.3.5} & 3.5 & \pageref{map:artigo_1:3.5} & Correspondência inspecionada\\[4pt]
\hyperref[map:artigo_1:3.6]{I.3.6} & 3.6 & \pageref{map:artigo_1:3.6} & Correspondência inspecionada\\[4pt]
\hyperref[map:artigo_1:4.1]{I.4.1} & 4.1 & \pageref{map:artigo_1:4.1} & Integrado à base. Identificação com os tipos concretos da v2 provada nos módulos NormalizationAtoms, Contracts, Enumerations, Classes, Errors e ReLU. Hipóteses, domínio, função e orçamentos preservados.\\[4pt]
\hyperref[map:artigo_1:4.2]{I.4.2} & 4.2 & \pageref{map:artigo_1:4.2} & Correspondência inspecionada\\[4pt]
\hyperref[map:artigo_1:4.3]{I.4.3} & 4.3 & \pageref{map:artigo_1:4.3} & Divergência documentada. O operador de extensão normalizado foi integrado à classe concreta da Fase 2. O enunciado TeX deste volume lê $g(t_0)$ como o valor em $t_0$ da primeira peça ($t_0$ pode estar fora de $K$); essa leitura não foi reverificada contra o enunciado Lean.\\[4pt]
\hyperref[map:artigo_1:4.4]{I.4.4} & 4.4 & \pageref{map:artigo_1:4.4} & Correspondência inspecionada\\[4pt]
\hyperref[map:artigo_1:4.5]{I.4.5} & 4.5 & \pageref{map:artigo_1:4.5} & Correspondência inspecionada\\[4pt]
\hyperref[map:artigo_1:5.1]{I.5.1} & 5.1 & \pageref{map:artigo_1:5.1} & Correspondência inspecionada\\[4pt]
\hyperref[map:artigo_1:5.2]{I.5.2} & 5.2 & \pageref{map:artigo_1:5.2} & Correspondência inspecionada\\[4pt]
\hyperref[map:artigo_1:6.1]{I.6.1} & 7.1 & \pageref{map:artigo_1:6.1} & Correspondência inspecionada\\[4pt]
\hyperref[map:artigo_1:6.2]{I.6.2} & 7.2 & \pageref{map:artigo_1:6.2} & Correspondência inspecionada\\[4pt]
\hyperref[map:artigo_1:6.3]{I.6.3} & 7.3 & \pageref{map:artigo_1:6.3} & Formalizado e compilado. Provas: \code{PendenciasParte1.certRefine, PendenciasParte1.certified_refinement} (raízes contidas em $[\alpha,\beta]$) e, para coberturas diádicas externas, \code{CoberturaExterna.certRefineProjected, CoberturaExterna.projected_certificate_sound}. Ver declarações e contratos nas fontes Lean. Compilado com Lean/Mathlib 4.28.0 (\texttt{lake build PendenciasParte1}); axiomas: \texttt{propext}, \texttt{Classical.choice}, \texttt{Quot.sound}. Ressalva: o refinamento é uma definição \texttt{noncomputable}, e as escolhas nodais e os endereços são dados por existência, como no texto.\\[4pt]
\hyperref[map:artigo_1:6.4]{I.6.4} & 7.4 & \pageref{map:artigo_1:6.4} & Formalizado e compilado. Provas: \code{PendenciasParte1.certified_refinement} (raízes contidas em $[\alpha,\beta]$) e, para coberturas diádicas externas, \code{CoberturaExterna.refinement_contract_external, CoberturaExterna.certified_refinement_external}, com o exemplo $[1/3,2/3]$ em \code{CoberturaExterna.thirds_success}. Ver declarações e contratos nas fontes Lean. Compilado com Lean/Mathlib 4.28.0 (\texttt{lake build PendenciasParte1}); axiomas: \texttt{propext}, \texttt{Classical.choice}, \texttt{Quot.sound}. Ressalva: término e correção sob o contrato de majorante (hipótese de entrada), sem certificar uma implementação executável ligada ao avaliador intervalar.\\[4pt]
\hyperref[map:artigo_1:6.5]{I.6.5} & 7.5 & \pageref{map:artigo_1:6.5} & Correspondência inspecionada\\[4pt]
\hyperref[map:artigo_1:6.6]{I.6.6} & 7.6 & \pageref{map:artigo_1:6.6} & Correspondência inspecionada\\[4pt]
\bottomrule\end{longtable}\normalsize

\subsection*{Parte II}
\small
\begin{longtable}{@{}p{.11\linewidth}p{.09\linewidth}p{.09\linewidth}p{.59\linewidth}@{}}
\toprule Resultado & Original & Página & Estado e ressalva\\\midrule\endhead
\hyperref[map:normalizacao:2.1]{II.2.1} & -- & \pageref{map:normalizacao:2.1} & Correspondência inspecionada. Resultado recebido preservado. A identificação dos tipos de átomos, endereços, classes e erros com a v2 está provada nas pontes Normalization; consultar INTEGRACAO\_NORMALIZACAO.md.\\[4pt]
\hyperref[map:normalizacao:2.2]{II.2.2} & -- & \pageref{map:normalizacao:2.2} & Correspondência inspecionada. Resultado recebido preservado. A identificação dos tipos de átomos, endereços, classes e erros com a v2 está provada nas pontes Normalization; consultar INTEGRACAO\_NORMALIZACAO.md.\\[4pt]
\hyperref[map:normalizacao:2.3]{II.2.3} & -- & \pageref{map:normalizacao:2.3} & Correspondência inspecionada. Resultado recebido preservado. A identificação dos tipos de átomos, endereços, classes e erros com a v2 está provada nas pontes Normalization; consultar INTEGRACAO\_NORMALIZACAO.md.\\[4pt]
\hyperref[map:normalizacao:2.4]{II.2.4} & -- & \pageref{map:normalizacao:2.4} & Integrado à base. Identificação com os tipos concretos da v2 provada nos módulos NormalizationAtoms, Contracts, Enumerations, Classes, Errors e ReLU. Hipóteses, domínio, função e orçamentos preservados.\\[4pt]
\hyperref[map:normalizacao:2.5]{II.2.5} & -- & \pageref{map:normalizacao:2.5} & Integrado à base. Identificação com os tipos concretos da v2 provada nos módulos NormalizationAtoms, Contracts, Enumerations, Classes, Errors e ReLU. Hipóteses, domínio, função e orçamentos preservados.\\[4pt]
\hyperref[map:normalizacao:3.1]{II.3.1} & -- & \pageref{map:normalizacao:3.1} & Integrado à base. Identificação com os tipos concretos da v2 provada nos módulos NormalizationAtoms, Contracts, Enumerations, Classes, Errors e ReLU. Hipóteses, domínio, função e orçamentos preservados.\\[4pt]
\hyperref[map:normalizacao:3.2]{II.3.2} & -- & \pageref{map:normalizacao:3.2} & Correspondência inspecionada. Resultado recebido preservado. A identificação dos tipos de átomos, endereços, classes e erros com a v2 está provada nas pontes Normalization; consultar INTEGRACAO\_NORMALIZACAO.md.\\[4pt]
\hyperref[map:normalizacao:3.3]{II.3.3} & -- & \pageref{map:normalizacao:3.3} & Correspondência inspecionada. Resultado recebido preservado. A identificação dos tipos de átomos, endereços, classes e erros com a v2 está provada nas pontes Normalization; consultar INTEGRACAO\_NORMALIZACAO.md.\\[4pt]
\hyperref[map:normalizacao:3.4]{II.3.4} & -- & \pageref{map:normalizacao:3.4} & Correspondência inspecionada. Resultado recebido preservado. A identificação dos tipos de átomos, endereços, classes e erros com a v2 está provada nas pontes Normalization; consultar INTEGRACAO\_NORMALIZACAO.md.\\[4pt]
\hyperref[map:normalizacao:3.5]{II.3.5} & -- & \pageref{map:normalizacao:3.5} & Integrado à base. Identificação com os tipos concretos da v2 provada nos módulos NormalizationAtoms, Contracts, Enumerations, Classes, Errors e ReLU. Hipóteses, domínio, função e orçamentos preservados.\\[4pt]
\hyperref[map:normalizacao:3.6]{II.3.6} & -- & \pageref{map:normalizacao:3.6} & Integrado à base. Identificação com os tipos concretos da v2 provada nos módulos NormalizationAtoms, Contracts, Enumerations, Classes, Errors e ReLU. Hipóteses, domínio, função e orçamentos preservados.\\[4pt]
\hyperref[map:normalizacao:4.1]{II.4.1} & -- & \pageref{map:normalizacao:4.1} & Correspondência inspecionada. Resultado recebido preservado. A identificação dos tipos de átomos, endereços, classes e erros com a v2 está provada nas pontes Normalization; consultar INTEGRACAO\_NORMALIZACAO.md.\\[4pt]
\hyperref[map:normalizacao:4.2]{II.4.2} & -- & \pageref{map:normalizacao:4.2} & Integrado à base. Identificação com os tipos concretos da v2 provada nos módulos NormalizationAtoms, Contracts, Enumerations, Classes, Errors e ReLU. Hipóteses, domínio, função e orçamentos preservados.\\[4pt]
\hyperref[map:normalizacao:4.3]{II.4.3} & -- & \pageref{map:normalizacao:4.3} & Integrado à base. Identificação com os tipos concretos da v2 provada nos módulos NormalizationAtoms, Contracts, Enumerations, Classes, Errors e ReLU. Hipóteses, domínio, função e orçamentos preservados.\\[4pt]
\hyperref[map:normalizacao:4.4]{II.4.4} & -- & \pageref{map:normalizacao:4.4} & Correspondência inspecionada. Resultado recebido preservado. A identificação dos tipos de átomos, endereços, classes e erros com a v2 está provada nas pontes Normalization; consultar INTEGRACAO\_NORMALIZACAO.md.\\[4pt]
\hyperref[map:normalizacao:4.5]{II.4.5} & -- & \pageref{map:normalizacao:4.5} & Integrado à base. Identificação com os tipos concretos da v2 provada nos módulos NormalizationAtoms, Contracts, Enumerations, Classes, Errors e ReLU. Hipóteses, domínio, função e orçamentos preservados.\\[4pt]
\hyperref[map:normalizacao:4.6]{II.4.6} & -- & \pageref{map:normalizacao:4.6} & Correspondência inspecionada. Resultado recebido preservado. A identificação dos tipos de átomos, endereços, classes e erros com a v2 está provada nas pontes Normalization; consultar INTEGRACAO\_NORMALIZACAO.md.\\[4pt]
\hyperref[map:normalizacao:4.7]{II.4.7} & -- & \pageref{map:normalizacao:4.7} & Integrado à base. Identificação com os tipos concretos da v2 provada nos módulos NormalizationAtoms, Contracts, Enumerations, Classes, Errors e ReLU. Hipóteses, domínio, função e orçamentos preservados.\\[4pt]
\hyperref[map:normalizacao:5.1]{II.5.1} & -- & \pageref{map:normalizacao:5.1} & Correspondência inspecionada. Resultado recebido preservado. A identificação dos tipos de átomos, endereços, classes e erros com a v2 está provada nas pontes Normalization; consultar INTEGRACAO\_NORMALIZACAO.md.\\[4pt]
\hyperref[map:normalizacao:5.2]{II.5.2} & -- & \pageref{map:normalizacao:5.2} & Integrado à base. Identificação com os tipos concretos da v2 provada nos módulos NormalizationAtoms, Contracts, Enumerations, Classes, Errors e ReLU. Hipóteses, domínio, função e orçamentos preservados.\\[4pt]
\hyperref[map:normalizacao:5.3]{II.5.3} & -- & \pageref{map:normalizacao:5.3} & Integrado à base. Identificação com os tipos concretos da v2 provada nos módulos NormalizationAtoms, Contracts, Enumerations, Classes, Errors e ReLU. Hipóteses, domínio, função e orçamentos preservados.\\[4pt]
\hyperref[map:normalizacao:5.4]{II.5.4} & -- & \pageref{map:normalizacao:5.4} & Correspondência inspecionada. Resultado recebido preservado. A identificação dos tipos de átomos, endereços, classes e erros com a v2 está provada nas pontes Normalization; consultar INTEGRACAO\_NORMALIZACAO.md.\\[4pt]
\hyperref[map:normalizacao:5.5]{II.5.5} & -- & \pageref{map:normalizacao:5.5} & Correspondência inspecionada. Resultado recebido preservado. A identificação dos tipos de átomos, endereços, classes e erros com a v2 está provada nas pontes Normalization; consultar INTEGRACAO\_NORMALIZACAO.md.\\[4pt]
\hyperref[map:normalizacao:5.6]{II.5.6} & -- & \pageref{map:normalizacao:5.6} & Integrado à base. Identificação com os tipos concretos da v2 provada nos módulos NormalizationAtoms, Contracts, Enumerations, Classes, Errors e ReLU. Hipóteses, domínio, função e orçamentos preservados.\\[4pt]
\hyperref[map:normalizacao:6.1]{II.6.1} & -- & \pageref{map:normalizacao:6.1} & Integrado à base. Identificação com os tipos concretos da v2 provada nos módulos NormalizationAtoms, Contracts, Enumerations, Classes, Errors e ReLU. Hipóteses, domínio, função e orçamentos preservados.\\[4pt]
\hyperref[map:normalizacao:6.2]{II.6.2} & -- & \pageref{map:normalizacao:6.2} & Integrado à base. Identificação com os tipos concretos da v2 provada nos módulos NormalizationAtoms, Contracts, Enumerations, Classes, Errors e ReLU. Hipóteses, domínio, função e orçamentos preservados.\\[4pt]
\hyperref[map:normalizacao:6.3]{II.6.3} & -- & \pageref{map:normalizacao:6.3} & Integrado à base. Identificação com os tipos concretos da v2 provada nos módulos NormalizationAtoms, Contracts, Enumerations, Classes, Errors e ReLU. Hipóteses, domínio, função e orçamentos preservados.\\[4pt]
\hyperref[map:normalizacao:6.4]{II.6.4} & -- & \pageref{map:normalizacao:6.4} & Integrado à base. Identificação com os tipos concretos da v2 provada nos módulos NormalizationAtoms, Contracts, Enumerations, Classes, Errors e ReLU. Hipóteses, domínio, função e orçamentos preservados.\\[4pt]
\hyperref[map:normalizacao:6.5]{II.6.5} & -- & \pageref{map:normalizacao:6.5} & Correspondência inspecionada. Resultado recebido preservado. A identificação dos tipos de átomos, endereços, classes e erros com a v2 está provada nas pontes Normalization; consultar INTEGRACAO\_NORMALIZACAO.md.\\[4pt]
\hyperref[map:normalizacao:7.1]{II.7.1} & -- & \pageref{map:normalizacao:7.1} & Correspondência inspecionada. Resultado recebido preservado. A identificação dos tipos de átomos, endereços, classes e erros com a v2 está provada nas pontes Normalization; consultar INTEGRACAO\_NORMALIZACAO.md.\\[4pt]
\hyperref[map:normalizacao:7.2]{II.7.2} & -- & \pageref{map:normalizacao:7.2} & Correspondência inspecionada. Resultado recebido preservado. A identificação dos tipos de átomos, endereços, classes e erros com a v2 está provada nas pontes Normalization; consultar INTEGRACAO\_NORMALIZACAO.md.\\[4pt]
\hyperref[map:normalizacao:7.3]{II.7.3} & -- & \pageref{map:normalizacao:7.3} & Correspondência inspecionada. Resultado recebido preservado. A identificação dos tipos de átomos, endereços, classes e erros com a v2 está provada nas pontes Normalization; consultar INTEGRACAO\_NORMALIZACAO.md.\\[4pt]
\hyperref[map:normalizacao:7.4]{II.7.4} & -- & \pageref{map:normalizacao:7.4} & Correspondência inspecionada. Resultado recebido preservado. A identificação dos tipos de átomos, endereços, classes e erros com a v2 está provada nas pontes Normalization; consultar INTEGRACAO\_NORMALIZACAO.md.\\[4pt]
\hyperref[map:normalizacao:8.1]{II.8.1} & -- & \pageref{map:normalizacao:8.1} & Correspondência inspecionada. Resultado recebido preservado. A identificação dos tipos de átomos, endereços, classes e erros com a v2 está provada nas pontes Normalization; consultar INTEGRACAO\_NORMALIZACAO.md.\\[4pt]
\hyperref[map:normalizacao:8.2]{II.8.2} & -- & \pageref{map:normalizacao:8.2} & Correspondência inspecionada. Resultado recebido preservado. A identificação dos tipos de átomos, endereços, classes e erros com a v2 está provada nas pontes Normalization; consultar INTEGRACAO\_NORMALIZACAO.md.\\[4pt]
\hyperref[map:normalizacao:8.3]{II.8.3} & -- & \pageref{map:normalizacao:8.3} & Correspondência inspecionada. Resultado recebido preservado. A identificação dos tipos de átomos, endereços, classes e erros com a v2 está provada nas pontes Normalization; consultar INTEGRACAO\_NORMALIZACAO.md.\\[4pt]
\hyperref[map:normalizacao:8.4]{II.8.4} & -- & \pageref{map:normalizacao:8.4} & Integrado à base. Identificação com os tipos concretos da v2 provada nos módulos NormalizationAtoms, Contracts, Enumerations, Classes, Errors e ReLU. Hipóteses, domínio, função e orçamentos preservados.\\[4pt]
\hyperref[map:normalizacao:8.5]{II.8.5} & -- & \pageref{map:normalizacao:8.5} & Correspondência inspecionada. Resultado recebido preservado. A identificação dos tipos de átomos, endereços, classes e erros com a v2 está provada nas pontes Normalization; consultar INTEGRACAO\_NORMALIZACAO.md.\\[4pt]
\hyperref[map:normalizacao:8.6]{II.8.6} & -- & \pageref{map:normalizacao:8.6} & Integrado à base. Identificação com os tipos concretos da v2 provada nos módulos NormalizationAtoms, Contracts, Enumerations, Classes, Errors e ReLU. Hipóteses, domínio, função e orçamentos preservados.\\[4pt]
\hyperref[map:normalizacao:8.7]{II.8.7} & -- & \pageref{map:normalizacao:8.7} & Integrado à base. Identificação com os tipos concretos da v2 provada nos módulos NormalizationAtoms, Contracts, Enumerations, Classes, Errors e ReLU. Hipóteses, domínio, função e orçamentos preservados.\\[4pt]
\bottomrule\end{longtable}\normalsize

\subsection*{Parte III}
\small
\begin{longtable}{@{}p{.11\linewidth}p{.09\linewidth}p{.09\linewidth}p{.59\linewidth}@{}}
\toprule Resultado & Original & Página & Estado e ressalva\\\midrule\endhead
\hyperref[map:artigo_2:2.1]{III.2.1} & -- & \pageref{map:artigo_2:2.1} & Correspondência inspecionada\\[4pt]
\hyperref[map:artigo_2:2.2]{III.2.2} & -- & \pageref{map:artigo_2:2.2} & Correspondência inspecionada\\[4pt]
\hyperref[map:artigo_2:2.3]{III.2.3} & -- & \pageref{map:artigo_2:2.3} & Formalizado e compilado. Cota de interpolação provada com \texttt{ContDiffOn} de ordem 2 em \texttt{Icc a b} e \texttt{derivWithin}, via extensão $C^2$ portada da Fase 3 (\texttt{Pendencias/Extensao.lean}); densidade e forma ReLU já batiam. Compilado com Lean/Mathlib 4.28.0 (\texttt{lake build Pendencias}); axiomas: \texttt{propext}, \texttt{Classical.choice}, \texttt{Quot.sound}. Ressalva anterior: Resultados globais ou classe FRM usam regularidade $C^2$ na reta. O artigo admite $C^2$ no intervalo; transportar a extensão Fase3.extF e preservar valores, erros e curvatura. Conferir individualmente as hipóteses dos componentes citados.\\[4pt]
\hyperref[map:artigo_2:3.1]{III.3.1} & 8.1 & \pageref{map:artigo_2:3.1} & Correspondência inspecionada\\[4pt]
\hyperref[map:artigo_2:3.2]{III.3.2} & 8.2 & \pageref{map:artigo_2:3.2} & Correspondência inspecionada\\[4pt]
\hyperref[map:artigo_2:3.3]{III.3.3} & 8.3 & \pageref{map:artigo_2:3.3} & Correspondência inspecionada\\[4pt]
\hyperref[map:artigo_2:3.4]{III.3.4} & 8.4 & \pageref{map:artigo_2:3.4} & Correspondência inspecionada\\[4pt]
\hyperref[map:artigo_2:3.5]{III.3.5} & 8.5 & \pageref{map:artigo_2:3.5} & Observação editorial. Observação ou diagnóstico do texto; não é um teorema Lean autônomo.\\[4pt]
\hyperref[map:artigo_2:3.6]{III.3.6} & 8.6 & \pageref{map:artigo_2:3.6} & Correspondência inspecionada\\[4pt]
\hyperref[map:artigo_2:3.7]{III.3.7} & 8.7 & \pageref{map:artigo_2:3.7} & Correspondência inspecionada\\[4pt]
\hyperref[map:artigo_2:3.8]{III.3.8} & 8.8 & \pageref{map:artigo_2:3.8} & Observação editorial. Observação ou diagnóstico do texto; não é um teorema Lean autônomo.\\[4pt]
\hyperref[map:artigo_2:3.9]{III.3.9} & 8.9 & \pageref{map:artigo_2:3.9} & Correspondência inspecionada\\[4pt]
\hyperref[map:artigo_2:3.10]{III.3.10} & 8.10 & \pageref{map:artigo_2:3.10} & Observação editorial. Observação ou diagnóstico do texto; não é um teorema Lean autônomo.\\[4pt]
\hyperref[map:artigo_2:3.11]{III.3.11} & 8.11 & \pageref{map:artigo_2:3.11} & Correspondência inspecionada\\[4pt]
\hyperref[map:artigo_2:3.12]{III.3.12} & 8.12 & \pageref{map:artigo_2:3.12} & Formalizado e compilado. Provas: \code{Pendencias.coding_formats_comparison}. Ver declarações e contratos nas fontes Lean. Compilado com Lean/Mathlib 4.28.0 (\texttt{lake build Pendencias}); axiomas: \texttt{propext}, \texttt{Classical.choice}, \texttt{Quot.sound}. Ressalva anterior, superada: \code{coding_gap_family} provava só a cadeia de átomos para $N$ par; a árvore de raiz única $J_{H-1-r,0}$, com $N=2^r$, é construída por \code{Pendencias.famTree} e tratada em \code{Pendencias.coding_formats_comparison}.\\[4pt]
\hyperref[map:artigo_2:3.13]{III.3.13} & 8.13 & \pageref{map:artigo_2:3.13} & Correspondência inspecionada\\[4pt]
\hyperref[map:artigo_2:3.14]{III.3.14} & 8.14 & \pageref{map:artigo_2:3.14} & Correspondência inspecionada\\[4pt]
\hyperref[map:artigo_2:3.15]{III.3.15} & 8.15 & \pageref{map:artigo_2:3.15} & Formalizado e compilado. Classe \texttt{FRMIcc R M} ($C^2$ apenas em $[0,1]$, $f''$ unilateral); $\mathrm{FRM}\subseteq\mathrm{FRMIcc}$ e \texttt{extF f 0 1} pertence a FRM e coincide com $f$ em $[0,1]$; $|f'|\le 4R+M/2$ com $f'$ unilateral. Compilado com Lean/Mathlib 4.28.0 (\texttt{lake build Pendencias}); axiomas: \texttt{propext}, \texttt{Classical.choice}, \texttt{Quot.sound}. Ressalva anterior: Resultados globais ou classe FRM usam regularidade $C^2$ na reta. O artigo admite $C^2$ no intervalo; transportar a extensão Fase3.extF e preservar valores, erros e curvatura. Conferir individualmente as hipóteses dos componentes citados.\\[4pt]
\hyperref[map:artigo_2:3.16]{III.3.16} & 8.16 & \pageref{map:artigo_2:3.16} & Formalizado e compilado. Classe \texttt{FRMIcc R M} como em III.3.15; mesma árvore determinística \texttt{uniformTree (jvals f r)}, que só amostra $f$ em $[0,1]$, e consequência para todo $N\ge 2$. Compilado com Lean/Mathlib 4.28.0 (\texttt{lake build Pendencias}); axiomas: \texttt{propext}, \texttt{Classical.choice}, \texttt{Quot.sound}. Ressalva anterior: Resultados globais ou classe FRM usam regularidade $C^2$ na reta. O artigo admite $C^2$ no intervalo; transportar a extensão Fase3.extF e preservar valores, erros e curvatura. Conferir individualmente as hipóteses dos componentes citados.\\[4pt]
\hyperref[map:artigo_2:3.17]{III.3.17} & 8.17 & \pageref{map:artigo_2:3.17} & Formalizado e compilado. Classe \texttt{FRMIcc R M} como em III.3.15; $n_\varepsilon$ e $H_\varepsilon$ inalterados; \texttt{jointErr}, \texttt{heightFor} e \texttt{supDist} só dependem de $f$ em $[0,1]$. Compilado com Lean/Mathlib 4.28.0 (\texttt{lake build Pendencias}); axiomas: \texttt{propext}, \texttt{Classical.choice}, \texttt{Quot.sound}. Ressalva anterior: Resultados globais ou classe FRM usam regularidade $C^2$ na reta. O artigo admite $C^2$ no intervalo; transportar a extensão Fase3.extF e preservar valores, erros e curvatura. Conferir individualmente as hipóteses dos componentes citados.\\[4pt]
\hyperref[map:artigo_2:3.18]{III.3.18} & 8.18 & \pageref{map:artigo_2:3.18} & Formalizado e compilado. Classe \texttt{FRMIcc R M} como em III.3.15; as coberturas de FRMIcc e FRM coincidem (\texttt{isCover\_iff}), logo as entropias coincidem; codebook transportado. Compilado com Lean/Mathlib 4.28.0 (\texttt{lake build Pendencias}); axiomas: \texttt{propext}, \texttt{Classical.choice}, \texttt{Quot.sound}. Ressalva anterior: Resultados globais ou classe FRM usam regularidade $C^2$ na reta. O artigo admite $C^2$ no intervalo; transportar a extensão Fase3.extF e preservar valores, erros e curvatura. Conferir individualmente as hipóteses dos componentes citados.\\[4pt]
\hyperref[map:artigo_2:3.19]{III.3.19} & 8.19 & \pageref{map:artigo_2:3.19} & Correspondência inspecionada\\[4pt]
\hyperref[map:artigo_2:3.20]{III.3.20} & 8.20 & \pageref{map:artigo_2:3.20} & Correspondência inspecionada\\[4pt]
\hyperref[map:artigo_2:3.21]{III.3.21} & 8.21 & \pageref{map:artigo_2:3.21} & Correspondência inspecionada\\[4pt]
\hyperref[map:artigo_2:3.22]{III.3.22} & 8.22 & \pageref{map:artigo_2:3.22} & Correspondência inspecionada\\[4pt]
\hyperref[map:artigo_2:3.23]{III.3.23} & 8.23 & \pageref{map:artigo_2:3.23} & Correspondência inspecionada\\[4pt]
\hyperref[map:artigo_2:3.24]{III.3.24} & 8.24 & \pageref{map:artigo_2:3.24} & Correspondência inspecionada\\[4pt]
\hyperref[map:artigo_2:3.25]{III.3.25} & 8.25 & \pageref{map:artigo_2:3.25} & Formalizado e compilado. Provas: \code{Pendencias.gen_eq_RepCat, Pendencias.enumSelect_eq_canon}. Ver declarações e contratos nas fontes Lean. Compilado com Lean/Mathlib 4.28.0 (\texttt{lake build Pendencias}); axiomas: \texttt{propext}, \texttt{Classical.choice}, \texttt{Quot.sound}. Ressalva anterior: A recorrência ou existência matemática tem prova; falta implementar e certificar a execução e a reconstrução com comparação exata.\\[4pt]
\hyperref[map:artigo_2:4.1]{III.4.1} & 9.1 & \pageref{map:artigo_2:4.1} & Correspondência inspecionada\\[4pt]
\hyperref[map:artigo_2:4.2]{III.4.2} & 9.2 & \pageref{map:artigo_2:4.2} & Correspondência inspecionada\\[4pt]
\hyperref[map:artigo_2:4.3]{III.4.3} & 9.3 & \pageref{map:artigo_2:4.3} & Correspondência inspecionada\\[4pt]
\hyperref[map:artigo_2:4.4]{III.4.4} & 9.4 & \pageref{map:artigo_2:4.4} & Correspondência inspecionada\\[4pt]
\hyperref[map:artigo_2:4.5]{III.4.5} & 9.5 & \pageref{map:artigo_2:4.5} & Correspondência inspecionada\\[4pt]
\hyperref[map:artigo_2:4.6]{III.4.6} & 9.6 & \pageref{map:artigo_2:4.6} & Observação editorial. Observação ou diagnóstico do texto; não é um teorema Lean autônomo.\\[4pt]
\hyperref[map:artigo_2:5.1]{III.5.1} & 10.1 & \pageref{map:artigo_2:5.1} & Correspondência inspecionada\\[4pt]
\hyperref[map:artigo_2:5.2]{III.5.2} & 10.2 & \pageref{map:artigo_2:5.2} & Correspondência inspecionada\\[4pt]
\hyperref[map:artigo_2:6.1]{III.6.1} & 11.1 & \pageref{map:artigo_2:6.1} & Correspondência inspecionada\\[4pt]
\hyperref[map:artigo_2:6.2]{III.6.2} & 11.2 & \pageref{map:artigo_2:6.2} & Correspondência inspecionada\\[4pt]
\hyperref[map:artigo_2:6.3]{III.6.3} & 11.3 & \pageref{map:artigo_2:6.3} & Formalizado e compilado. Provas: \code{Pendencias.DPprog.fixTable_cost, Pendencias.DPprog.fixRoot_spec}. Ver declarações e contratos nas fontes Lean. Compilado com Lean/Mathlib 4.28.0 (\texttt{lake build Pendencias}); axiomas: \texttt{propext}, \texttt{Classical.choice}, \texttt{Quot.sound}. Ressalva anterior: A recorrência ou existência matemática tem prova; falta implementar e certificar a execução e a reconstrução com comparação exata.\\[4pt]
\hyperref[map:artigo_2:7.1]{III.7.1} & 12.1 & \pageref{map:artigo_2:7.1} & Correspondência inspecionada\\[4pt]
\hyperref[map:artigo_2:7.2]{III.7.2} & 12.2 & \pageref{map:artigo_2:7.2} & Correspondência inspecionada\\[4pt]
\hyperref[map:artigo_2:7.3]{III.7.3} & 12.3 & \pageref{map:artigo_2:7.3} & Correspondência inspecionada\\[4pt]
\hyperref[map:artigo_2:7.4]{III.7.4} & 12.4 & \pageref{map:artigo_2:7.4} & Formalizado e compilado. Provas: \code{Pendencias.DPprog.varTable_cost, Pendencias.DPprog.varRoot_spec}. Ver declarações e contratos nas fontes Lean. Compilado com Lean/Mathlib 4.28.0 (\texttt{lake build Pendencias}); axiomas: \texttt{propext}, \texttt{Classical.choice}, \texttt{Quot.sound}. Ressalva anterior: A recorrência ou existência matemática tem prova; falta implementar e certificar a execução e a reconstrução com comparação exata.\\[4pt]
\hyperref[map:artigo_2:8.1]{III.8.1} & 13.1 & \pageref{map:artigo_2:8.1} & Correspondência inspecionada\\[4pt]
\hyperref[map:artigo_2:8.2]{III.8.2} & 13.2 & \pageref{map:artigo_2:8.2} & Correspondência inspecionada\\[4pt]
\hyperref[map:artigo_2:8.3]{III.8.3} & 13.3 & \pageref{map:artigo_2:8.3} & Observação editorial. Observação ou diagnóstico do texto; não é um teorema Lean autônomo.\\[4pt]
\hyperref[map:artigo_2:8.4]{III.8.4} & 13.4 & \pageref{map:artigo_2:8.4} & Correspondência inspecionada\\[4pt]
\hyperref[map:artigo_2:9.1]{III.9.1} & 14.1 & \pageref{map:artigo_2:9.1} & Formalizado e compilado. Complexidade em número de estados: prova \code{DenseDP.state_complexity_dense}; modelo de custo unitário explícito (leitura/escrita da tabela, soma, comparação); constantes 8 e 7. Compilado com Lean/Mathlib 4.28.0 (\texttt{lake build Pendencias}); axiomas: \texttt{propext}, \texttt{Classical.choice}, \texttt{Quot.sound}. Ressalva anterior: Prova cardinalidades das tabelas e transições; não certifica tempo e memória de um programa com execução ligada a essas cardinalidades.\\[4pt]
\hyperref[map:artigo_2:9.2]{III.9.2} & 14.2 & \pageref{map:artigo_2:9.2} & Observação editorial. Observação ou diagnóstico do texto; não é um teorema Lean autônomo.\\[4pt]
\hyperref[map:artigo_2:9.3]{III.9.3} & 14.3 & \pageref{map:artigo_2:9.3} & Formalizado e compilado. Complexidade de bits: prova \code{DenseDP.bit_complexity_dense}; custo de operação aritmética em bits é hipótese do modelo (\code{BitModel}). Compilado com Lean/Mathlib 4.28.0 (\texttt{lake build Pendencias}); axiomas: \texttt{propext}, \texttt{Classical.choice}, \texttt{Quot.sound}. Ressalva anterior: bit\_complexity limita o tamanho dos custos racionais. As cotas de tempo e memória em bits do artigo exigem semântica de execução e custo de operações/desempate.\\[4pt]
\hyperref[map:artigo_2:9.4]{III.9.4} & 14.4 & \pageref{map:artigo_2:9.4} & Correspondência inspecionada\\[4pt]
\hyperref[map:artigo_2:9.5]{III.9.5} & 14.5 & \pageref{map:artigo_2:9.5} & Correspondência inspecionada\\[4pt]
\hyperref[map:artigo_2:11.1]{III.11.1} & 16.1 & \pageref{map:artigo_2:11.1} & Correspondência inspecionada\\[4pt]
\hyperref[map:artigo_2:11.2]{III.11.2} & 16.2 & \pageref{map:artigo_2:11.2} & Formalizado e compilado. $f\in C^2(K)$ via \texttt{ContDiffOn}; $e_f(J)$, $f''$ e $\omega_{f''}$ unilaterais em $K$, por \texttt{extF} e lemas de congruência. Compilado com Lean/Mathlib 4.28.0 (\texttt{lake build Pendencias}); axiomas: \texttt{propext}, \texttt{Classical.choice}, \texttt{Quot.sound}. Ressalva anterior: Resultados globais ou classe FRM usam regularidade $C^2$ na reta. O artigo admite $C^2$ no intervalo; transportar a extensão Fase3.extF e preservar valores, erros e curvatura. Conferir individualmente as hipóteses dos componentes citados.\\[4pt]
\hyperref[map:artigo_2:11.3]{III.11.3} & 16.3 & \pageref{map:artigo_2:11.3} & Formalizado e compilado. $f\in C^2(K)$ via \texttt{ContDiffOn}, com \texttt{curvIOn} ($I(f)$ com $f''$ unilateral); as três formas (eventual, $\liminf/\limsup$, caso $I(f)=0$). Compilado com Lean/Mathlib 4.28.0 (\texttt{lake build Pendencias}); axiomas: \texttt{propext}, \texttt{Classical.choice}, \texttt{Quot.sound}. Ressalva anterior: Resultados globais ou classe FRM usam regularidade $C^2$ na reta. O artigo admite $C^2$ no intervalo; transportar a extensão Fase3.extF e preservar valores, erros e curvatura. Conferir individualmente as hipóteses dos componentes citados.\\[4pt]
\hyperref[map:artigo_2:12.1]{III.12.1} & 17.1 & \pageref{map:artigo_2:12.1} & Formalizado e compilado. Mesmas seis conclusões de \texttt{quadratic\_dyadic\_constant} com $q\in C^2(K)$ (\texttt{ContDiffOn}) e $q''\equiv\kappa$ unilateral em $K$; usa \texttt{extF} e lemas de congruência (\texttt{Edia\_congr}, \texttt{nThr\_congr}). Compilado com Lean/Mathlib 4.28.0 (\texttt{lake build Pendencias}); axiomas: \texttt{propext}, \texttt{Classical.choice}, \texttt{Quot.sound}. Ressalva anterior: Resultados globais ou classe FRM usam regularidade $C^2$ na reta. O artigo admite $C^2$ no intervalo; transportar a extensão Fase3.extF e preservar valores, erros e curvatura. Conferir individualmente as hipóteses dos componentes citados.\\[4pt]
\hyperref[map:artigo_2:12.2]{III.12.2} & 17.2 & \pageref{map:artigo_2:12.2} & Generalização documentada. A hipótese $C^1$ foi enfraquecida no texto para continuidade, que é o que a prova usa; o Lean tem essa mesma hipótese, pois a igualdade com quadráticas nas células fechadas conserva compatibilidade nos extremos.\\[4pt]
\hyperref[map:artigo_2:12.3]{III.12.3} & 17.3 & \pageref{map:artigo_2:12.3} & Prova textual sincronizada. Demonstração textual sincronizada com \code{Artigo.synchronized_phases}; usa contagem exata e periodicidade, sem o envelope $C^2$ global.\\[4pt]
\hyperref[map:artigo_2:12.4]{III.12.4} & 17.4 & \pageref{map:artigo_2:12.4} & Observação editorial. Observação ou diagnóstico do texto; não é um teorema Lean autônomo.\\[4pt]
\bottomrule\end{longtable}\normalsize

\subsection*{Parte IV}
\small
\begin{longtable}{@{}p{.11\linewidth}p{.09\linewidth}p{.09\linewidth}p{.59\linewidth}@{}}
\toprule Resultado & Original & Página & Estado e ressalva\\\midrule\endhead
\hyperref[map:artigo_3:3.1]{IV.3.1} & 18.1 & \pageref{map:artigo_3:3.1} & Formalizado e compilado. Definições \texttt{Sigma} e \texttt{pErr} sobre $K=[0,1]$ (domínio de \texttt{jointErr}); inclusão, monotonia, cota por $\|f\|$ e atingimento do ínfimo (catálogo finito), além da fórmula $\min\{\|f\|,E_{N,H(N)}\}$. Compilado com Lean/Mathlib 4.28.0 (\texttt{lake build Pendencias}); axiomas: \texttt{propext}, \texttt{Classical.choice}, \texttt{Quot.sound}. Ressalva anterior: Não foi localizado resultado Lean que certifique este enunciado. Fase3 é uma etapa da formalização do Artigo 1, não o atual Artigo 3.\\[4pt]
\hyperref[map:artigo_3:3.2]{IV.3.2} & 18.2 & \pageref{map:artigo_3:3.2} & Formalizado e compilado. Provas: \code{Pendencias.Agenda.dyadic_reduction, Pendencias.Agenda.dyadic_reduction_inf}. Ver declarações e contratos nas fontes Lean. Compilado com Lean/Mathlib 4.28.0 (\texttt{lake build Pendencias}); axiomas: \texttt{propext}, \texttt{Classical.choice}, \texttt{Quot.sound}. Ressalva anterior: Não foi localizado resultado Lean que certifique este enunciado. Fase3 é uma etapa da formalização do Artigo 1, não o atual Artigo 3.\\[4pt]
\hyperref[map:artigo_3:3.3]{IV.3.3} & 18.3 & \pageref{map:artigo_3:3.3} & Formalizado e compilado. Provas: \code{Pendencias.Agenda.homogeneity_failure, Pendencias.Agenda.homogeneity_failure_inf}. Ver declarações e contratos nas fontes Lean. Compilado com Lean/Mathlib 4.28.0 (\texttt{lake build Pendencias}); axiomas: \texttt{propext}, \texttt{Classical.choice}, \texttt{Quot.sound}. Ressalva anterior: Não foi localizado resultado Lean que certifique este enunciado. Fase3 é uma etapa da formalização do Artigo 1, não o atual Artigo 3.\\[4pt]
\hyperref[map:artigo_3:3.4]{IV.3.4} & 18.4 & \pageref{map:artigo_3:3.4} & Observação editorial. Não foi localizado resultado Lean que certifique este enunciado. Fase3 é uma etapa da formalização do Artigo 1, não o atual Artigo 3.\\[4pt]
\hyperref[map:artigo_3:4.1]{IV.4.1} & 19.1 & \pageref{map:artigo_3:4.1} & Formalizado e compilado. Jackson racional em raiz diádica geral, constante absoluta 8 e versão uniforme sobre bolas; Whitney demonstrada, sem hipótese adicional não provada. Compilado com Lean/Mathlib 4.28.0 (\texttt{lake build Pendencias}); axiomas: \texttt{propext}, \texttt{Classical.choice}, \texttt{Quot.sound}. Ressalva anterior: Não foi localizado resultado Lean que certifique este enunciado. Fase3 é uma etapa da formalização do Artigo 1, não o atual Artigo 3.\\[4pt]
\hyperref[map:artigo_3:4.2]{IV.4.2} & 19.2 & \pageref{map:artigo_3:4.2} & Formalizado e compilado. Inclusões diretas de Besov para os regimes do enunciado, inclusive agenda com folga polinomial, em raiz diádica geral. Compilado com Lean/Mathlib 4.28.0 (\texttt{lake build Pendencias}); axiomas: \texttt{propext}, \texttt{Classical.choice}, \texttt{Quot.sound}. Ressalva anterior: Não foi localizado resultado Lean que certifique este enunciado. Fase3 é uma etapa da formalização do Artigo 1, não o atual Artigo 3.\\[4pt]
\hyperref[map:artigo_3:4.3]{IV.4.3} & 19.3 & \pageref{map:artigo_3:4.3} & Formalizado e compilado. Cotas quantitativas uniformes sobre bolas nos três regimes, com constantes independentes da função e do raio. Compilado com Lean/Mathlib 4.28.0 (\texttt{lake build Pendencias}); axiomas: \texttt{propext}, \texttt{Classical.choice}, \texttt{Quot.sound}. Ressalva anterior: Não foi localizado resultado Lean que certifique este enunciado. Fase3 é uma etapa da formalização do Artigo 1, não o atual Artigo 3.\\[4pt]
\hyperref[map:artigo_3:4.4]{IV.4.4} & 19.4 & \pageref{map:artigo_3:4.4} & Observação editorial. Não foi localizado resultado Lean que certifique este enunciado. Fase3 é uma etapa da formalização do Artigo 1, não o atual Artigo 3.\\[4pt]
\hyperref[map:artigo_3:4.5]{IV.4.5} & 19.5 & \pageref{map:artigo_3:4.5} & Formalizado e compilado. Funcional \texttt{rateInf} ($\mathfrak{a}_{\alpha,\infty}$) e agenda \texttt{H2 R M N} $=\lceil C_H N^2\rceil$; supremo sobre \texttt{FRMIcc R M} de $\mathfrak{a}_{2,\infty}\le\max\{R,(M+4)/2\}$; para $N\ge 2$, montagem admissível com códigos $\le C_b N\log_2(2N)$. Base: \texttt{joint\_all\_N\_Icc}, \texttt{joint\_budget\_Icc}, \texttt{pErr\_eq\_min}. Compilado com Lean/Mathlib 4.28.0 (\texttt{lake build Pendencias}); axiomas: \texttt{propext}, \texttt{Classical.choice}, \texttt{Quot.sound}. Ressalva anterior: A estimativa $C^2$ de base existe. Falta transportá-la para as definições de agenda, erro e classe do atual Artigo 3; não contado como teorema integralmente formalizado.\\[4pt]
\hyperref[map:artigo_3:4.6]{IV.4.6} & 19.6 & \pageref{map:artigo_3:4.6} & Formalizado e compilado. Piso $e_N\ge 1/(2H(N))$, $\mathrm{dist}(1/(2H),\mathcal{Q}_H)=1/(2H)$ e a consequência $H(N)\ge (2C)^{-1}N^s$ (para $N\ge 1$), sobre $K=[0,1]$. Compilado com Lean/Mathlib 4.28.0 (\texttt{lake build Pendencias}); axiomas: \texttt{propext}, \texttt{Classical.choice}, \texttt{Quot.sound}. Ressalva anterior: Não foi localizado resultado Lean que certifique este enunciado. Fase3 é uma etapa da formalização do Artigo 1, não o atual Artigo 3.\\[4pt]
\hyperref[map:artigo_3:5.1]{IV.5.1} & 20.1 & \pageref{map:artigo_3:5.1} & Formalizado e compilado. Provas: \code{Pendencias.Marcada.marked_extends, Pendencias.Marcada.Atom.toM_injective}. Ver declarações e contratos nas fontes Lean. Compilado com Lean/Mathlib 4.28.0 (\texttt{lake build Pendencias}); axiomas: \texttt{propext}, \texttt{Classical.choice}, \texttt{Quot.sound}. Ressalva anterior: Não foi localizado resultado Lean que certifique este enunciado. Fase3 é uma etapa da formalização do Artigo 1, não o atual Artigo 3.\\[4pt]
\hyperref[map:artigo_3:5.2]{IV.5.2} & 20.2 & \pageref{map:artigo_3:5.2} & Formalizado e compilado. Reversibilidade rasa completa para parâmetros reais: equivalência aritmética, árvore canônica única (inclusive rótulos), cotas explícitas de folhas e condição de inclinações não nulas para a biblioteca original. Prova: \code{RealReLU.exact_shallow_reversal}; funções em $[0,1]$, incluindo $M=0$. Compilado com Lean/Mathlib 4.28.0 (\texttt{lake build Pendencias}); axiomas: \texttt{propext}, \texttt{Classical.choice}, \texttt{Quot.sound}. Ressalva anterior: Não foi localizado resultado Lean que certifique este enunciado. Fase3 é uma etapa da formalização do Artigo 1, não o atual Artigo 3.\\[4pt]
\hyperref[map:artigo_3:5.3]{IV.5.3} & 20.3 & \pageref{map:artigo_3:5.3} & Formalizado e compilado. Provas: \code{Pendencias.Marcada.RedReLU.Hcan, Pendencias.Marcada.RedReLU.Hcan_eq_maxHeight, Pendencias.Marcada.canonShape_no_interior, Pendencias.Marcada.RedReLU.canon_leaves_free}. Ver declarações e contratos nas fontes Lean. Compilado com Lean/Mathlib 4.28.0 (\texttt{lake build Pendencias}); axiomas: \texttt{propext}, \texttt{Classical.choice}, \texttt{Quot.sound}. Ressalva anterior: Não foi localizado resultado Lean que certifique este enunciado. Fase3 é uma etapa da formalização do Artigo 1, não o atual Artigo 3.\\[4pt]
\hyperref[map:artigo_3:5.4]{IV.5.4} & 20.4 & \pageref{map:artigo_3:5.4} & Formalizado e compilado. Formalizado nas duas direções; a cota inversa foi corrigida (fator $M+1$, constante $C=5$) porque a impressa é falsa para $M=0$. Provas: \code{Pendencias.Marcada.AlturaInversa.reverse_height_bit_bound} e \code{Pendencias.Marcada.Ancorada.direct_part}. Compilado com Lean/Mathlib 4.28.0 (\texttt{lake build Pendencias}); axiomas: \texttt{propext}, \texttt{Classical.choice}, \texttt{Quot.sound}. Ressalva anterior: Não foi localizado resultado Lean que certifique este enunciado. Fase3 é uma etapa da formalização do Artigo 1, não o atual Artigo 3.\\[4pt]
\hyperref[map:artigo_3:5.5]{IV.5.5} & 20.5 & \pageref{map:artigo_3:5.5} & Formalizado e compilado. Reversão aproximada quantizada com $C=320$. Prova: \code{QuantIV55.quantized_approximate_reverse}; erro em $[0,1]$, folhas da montagem e altura atômica máxima. Compilado com Lean/Mathlib 4.28.0 (\texttt{lake build Pendencias}); axiomas: \texttt{propext}, \texttt{Classical.choice}, \texttt{Quot.sound}. Ressalva anterior: Não foi localizado resultado Lean que certifique este enunciado. Fase3 é uma etapa da formalização do Artigo 1, não o atual Artigo 3.\\[4pt]
\hyperref[map:artigo_3:5.6]{IV.5.6} & 20.6 & \pageref{map:artigo_3:5.6} & Formalizado e compilado. Corolário para redes profundas: prova em \code{DeepIV56.deep_functional_reverse}; $R(F)$ é o mínimo de peças afins de $[0,1]$, $B$ e $L_F$ medidos em $[0,1]$. Compilado com Lean/Mathlib 4.28.0 (\texttt{lake build Pendencias}); axiomas: \texttt{propext}, \texttt{Classical.choice}, \texttt{Quot.sound}. Ressalva anterior: Não foi localizado resultado Lean que certifique este enunciado. Fase3 é uma etapa da formalização do Artigo 1, não o atual Artigo 3.\\[4pt]
\hyperref[map:artigo_3:5.7]{IV.5.7} & 20.7 & \pageref{map:artigo_3:5.7} & Formalizado e compilado. Provas: \code{Pendencias.Agenda.inverse_modulus, Pendencias.Agenda.inverse_besov}. Ver declarações e contratos nas fontes Lean. Compilado com Lean/Mathlib 4.28.0 (\texttt{lake build Pendencias}); axiomas: \texttt{propext}, \texttt{Classical.choice}, \texttt{Quot.sound}. Ressalva anterior: Não foi localizado resultado Lean que certifique este enunciado. Fase3 é uma etapa da formalização do Artigo 1, não o atual Artigo 3.\\[4pt]
\hyperref[map:artigo_3:5.8]{IV.5.8} & 20.8 & \pageref{map:artigo_3:5.8} & Formalizado e compilado. Provas: \code{Pendencias.Agenda.raw_library_inverse_inf, Pendencias.Agenda.raw_library_inverse, Pendencias.Agenda.schedule_raw_inverse, Pendencias.Agenda.vartheta_eq}. Ver declarações e contratos nas fontes Lean. Compilado com Lean/Mathlib 4.28.0 (\texttt{lake build Pendencias}); axiomas: \texttt{propext}, \texttt{Classical.choice}, \texttt{Quot.sound}. Ressalva anterior: Não foi localizado resultado Lean que certifique este enunciado. Fase3 é uma etapa da formalização do Artigo 1, não o atual Artigo 3.\\[4pt]
\hyperref[map:artigo_3:5.9]{IV.5.9} & 20.9 & \pageref{map:artigo_3:5.9} & Formalizado e compilado. Teorema formalizado na biblioteca original: continuidade, taxa para todo $N\ge1$ e não pertencimento às classes Besov prescritas. Prova: \code{Pendencias.PenteGemini.critical_tent_comb}. Compilado com Lean/Mathlib 4.28.0 (\texttt{lake build Pendencias}); axiomas: \texttt{propext}, \texttt{Classical.choice}, \texttt{Quot.sound}. Ressalva anterior: Não foi localizado resultado Lean que certifique este enunciado. Fase3 é uma etapa da formalização do Artigo 1, não o atual Artigo 3.\\[4pt]
\hyperref[map:artigo_3:5.10]{IV.5.10} & 20.10 & \pageref{map:artigo_3:5.10} & Formalizado e compilado. Provas: \code{Pendencias.Agenda.universal_inverse}. Ver declarações e contratos nas fontes Lean. Compilado com Lean/Mathlib 4.28.0 (\texttt{lake build Pendencias}); axiomas: \texttt{propext}, \texttt{Classical.choice}, \texttt{Quot.sound}. Ressalva anterior: Não foi localizado resultado Lean que certifique este enunciado. Fase3 é uma etapa da formalização do Artigo 1, não o atual Artigo 3.\\[4pt]
\hyperref[map:artigo_3:5.11]{IV.5.11} & 20.11 & \pageref{map:artigo_3:5.11} & Observação editorial. Não foi localizado resultado Lean que certifique este enunciado. Fase3 é uma etapa da formalização do Artigo 1, não o atual Artigo 3.\\[4pt]
\hyperref[map:artigo_3:5.12]{IV.5.12} & 20.12 & \pageref{map:artigo_3:5.12} & Formalizado e compilado. Jackson para o funcional de realização em $[0,1]$, cota bilateral e versão uniforme; constantes universais explícitas. Compilado com Lean/Mathlib 4.28.0 (\texttt{lake build Pendencias}); axiomas: \texttt{propext}, \texttt{Classical.choice}, \texttt{Quot.sound}. Ressalva anterior: Não foi localizado resultado Lean que certifique este enunciado. Fase3 é uma etapa da formalização do Artigo 1, não o atual Artigo 3.\\[4pt]
\hyperref[map:artigo_3:5.13]{IV.5.13} & 20.13 & \pageref{map:artigo_3:5.13} & Formalizado e compilado. Igualdades de classes de Besov e realização em todos os regimes declarados, incluindo folga polinomial. Compilado com Lean/Mathlib 4.28.0 (\texttt{lake build Pendencias}); axiomas: \texttt{propext}, \texttt{Classical.choice}, \texttt{Quot.sound}. Ressalva anterior: Não foi localizado resultado Lean que certifique este enunciado. Fase3 é uma etapa da formalização do Artigo 1, não o atual Artigo 3.\\[4pt]
\hyperref[map:artigo_3:5.14]{IV.5.14} & 20.14 & \pageref{map:artigo_3:5.14} & Formalizado e compilado. Estabilidade das igualdades para famílias aninhadas que contêm zero e as montagens na escala linear declarada. Compilado com Lean/Mathlib 4.28.0 (\texttt{lake build Pendencias}); axiomas: \texttt{propext}, \texttt{Classical.choice}, \texttt{Quot.sound}. Ressalva anterior: Não foi localizado resultado Lean que certifique este enunciado. Fase3 é uma etapa da formalização do Artigo 1, não o atual Artigo 3.\\[4pt]
\bottomrule\end{longtable}\normalsize

\subsection*{Parte V}
\small
\begin{longtable}{@{}p{.11\linewidth}p{.09\linewidth}p{.09\linewidth}p{.59\linewidth}@{}}
\toprule Resultado & Original & Página & Estado e ressalva\\\midrule\endhead
\hyperref[map:complemento_original:6.1]{V.6.1} & 6.1 & \pageref{map:complemento_original:6.1} & Formalizado e compilado. Provas: \code{Pendencias.tree_rate_pow, Pendencias.tree_rate, Pendencias.tree_rationalization}. Ver declarações e contratos nas fontes Lean. Compilado com Lean/Mathlib 4.28.0 (\texttt{lake build Pendencias}); axiomas: \texttt{propext}, \texttt{Classical.choice}, \texttt{Quot.sound}. Ressalva anterior: Resultado original integralmente preservado no complemento; a correspondência editorial não afirma prova Lean.\\[4pt]
\hyperref[map:complemento_original:6.2]{V.6.2} & 6.2 & \pageref{map:complemento_original:6.2} & Formalizado e compilado. Cota assintótica com nós livres e derivadas unilaterais para funções $C^2$ no intervalo fechado. Compilado com Lean/Mathlib 4.28.0 (\texttt{lake build Pendencias}); axiomas: \texttt{propext}, \texttt{Classical.choice}, \texttt{Quot.sound}. Ressalva anterior: Resultado original integralmente preservado no complemento; a correspondência editorial não afirma prova Lean.\\[4pt]
\hyperref[map:complemento_original:6.3]{V.6.3} & 6.3 & \pageref{map:complemento_original:6.3} & Formalizado e compilado. Estimativa direta em raiz diádica geral, constante de Whitney explícita e montagem racional com erro adicional arbitrário. Compilado com Lean/Mathlib 4.28.0 (\texttt{lake build Pendencias}); axiomas: \texttt{propext}, \texttt{Classical.choice}, \texttt{Quot.sound}. Ressalva anterior: Resultado original integralmente preservado no complemento; a correspondência editorial não afirma prova Lean.\\[4pt]
\hyperref[map:complemento_original:21.1]{V.21.1} & 21.1 & \pageref{map:complemento_original:21.1} & Formalizado e compilado. Provado; a forma \texttt{runge\_majorant} dispensa $J\subseteq[-1,1]$. Inclui \texttt{deriv2\_runge}. Compilado com Lean/Mathlib 4.28.0 (\texttt{lake build Pendencias}); axiomas: \texttt{propext}, \texttt{Classical.choice}, \texttt{Quot.sound}. Ressalva anterior: Resultado original integralmente preservado no complemento; a correspondência editorial não afirma prova Lean.\\[4pt]
\bottomrule\end{longtable}\normalsize
